% Options for packages loaded elsewhere
% Options for packages loaded elsewhere
\PassOptionsToPackage{unicode}{hyperref}
\PassOptionsToPackage{hyphens}{url}
%
\documentclass[
  ignorenonframetext,
]{beamer}
\newif\ifbibliography
\usepackage{pgfpages}
\setbeamertemplate{caption}[numbered]
\setbeamertemplate{caption label separator}{: }
\setbeamercolor{caption name}{fg=normal text.fg}
\beamertemplatenavigationsymbolsempty
% remove section numbering
\setbeamertemplate{part page}{
  \centering
  \begin{beamercolorbox}[sep=16pt,center]{part title}
    \usebeamerfont{part title}\insertpart\par
  \end{beamercolorbox}
}
\setbeamertemplate{section page}{
  \centering
  \begin{beamercolorbox}[sep=12pt,center]{section title}
    \usebeamerfont{section title}\insertsection\par
  \end{beamercolorbox}
}
\setbeamertemplate{subsection page}{
  \centering
  \begin{beamercolorbox}[sep=8pt,center]{subsection title}
    \usebeamerfont{subsection title}\insertsubsection\par
  \end{beamercolorbox}
}
% Prevent slide breaks in the middle of a paragraph
\widowpenalties 1 10000
\raggedbottom
\AtBeginPart{
  \frame{\partpage}
}
\AtBeginSection{
  \ifbibliography
  \else
    \frame{\sectionpage}
  \fi
}
\AtBeginSubsection{
  \frame{\subsectionpage}
}
\usepackage{iftex}
\ifPDFTeX
  \usepackage[T1]{fontenc}
  \usepackage[utf8]{inputenc}
  \usepackage{textcomp} % provide euro and other symbols
\else % if luatex or xetex
  \usepackage{unicode-math} % this also loads fontspec
  \defaultfontfeatures{Scale=MatchLowercase}
  \defaultfontfeatures[\rmfamily]{Ligatures=TeX,Scale=1}
\fi
\usepackage{lmodern}

\usetheme[]{Madrid}
\usecolortheme[]{dolphin}
\ifPDFTeX\else
  % xetex/luatex font selection
\fi
% Use upquote if available, for straight quotes in verbatim environments
\IfFileExists{upquote.sty}{\usepackage{upquote}}{}
\IfFileExists{microtype.sty}{% use microtype if available
  \usepackage[]{microtype}
  \UseMicrotypeSet[protrusion]{basicmath} % disable protrusion for tt fonts
}{}
\makeatletter
\@ifundefined{KOMAClassName}{% if non-KOMA class
  \IfFileExists{parskip.sty}{%
    \usepackage{parskip}
  }{% else
    \setlength{\parindent}{0pt}
    \setlength{\parskip}{6pt plus 2pt minus 1pt}}
}{% if KOMA class
  \KOMAoptions{parskip=half}}
\makeatother


\usepackage{longtable,booktabs,array}
\newcounter{none} % for unnumbered tables
\usepackage{calc} % for calculating minipage widths
\usepackage{caption}
% Make caption package work with longtable
\makeatletter
\def\fnum@table{\tablename~\thetable}
\makeatother
\usepackage{graphicx}
\makeatletter
\newsavebox\pandoc@box
\newcommand*\pandocbounded[1]{% scales image to fit in text height/width
  \sbox\pandoc@box{#1}%
  \Gscale@div\@tempa{\textheight}{\dimexpr\ht\pandoc@box+\dp\pandoc@box\relax}%
  \Gscale@div\@tempb{\linewidth}{\wd\pandoc@box}%
  \ifdim\@tempb\p@<\@tempa\p@\let\@tempa\@tempb\fi% select the smaller of both
  \ifdim\@tempa\p@<\p@\scalebox{\@tempa}{\usebox\pandoc@box}%
  \else\usebox{\pandoc@box}%
  \fi%
}
% Set default figure placement to htbp
\def\fps@figure{htbp}
\makeatother


% definitions for citeproc citations
\NewDocumentCommand\citeproctext{}{}
\NewDocumentCommand\citeproc{mm}{%
  \begingroup\def\citeproctext{#2}\cite{#1}\endgroup}
\makeatletter
 % allow citations to break across lines
 \let\@cite@ofmt\@firstofone
 % avoid brackets around text for \cite:
 \def\@biblabel#1{}
 \def\@cite#1#2{{#1\if@tempswa , #2\fi}}
\makeatother
\newlength{\cslhangindent}
\setlength{\cslhangindent}{1.5em}
\newlength{\csllabelwidth}
\setlength{\csllabelwidth}{3em}
\newenvironment{CSLReferences}[2] % #1 hanging-indent, #2 entry-spacing
 {\begin{list}{}{%
  \setlength{\itemindent}{0pt}
  \setlength{\leftmargin}{0pt}
  \setlength{\parsep}{0pt}
  % turn on hanging indent if param 1 is 1
  \ifodd #1
   \setlength{\leftmargin}{\cslhangindent}
   \setlength{\itemindent}{-1\cslhangindent}
  \fi
  % set entry spacing
  \setlength{\itemsep}{#2\baselineskip}}}
 {\end{list}}
\usepackage{calc}
\newcommand{\CSLBlock}[1]{\hfill\break\parbox[t]{\linewidth}{\strut\ignorespaces#1\strut}}
\newcommand{\CSLLeftMargin}[1]{\parbox[t]{\csllabelwidth}{\strut#1\strut}}
\newcommand{\CSLRightInline}[1]{\parbox[t]{\linewidth - \csllabelwidth}{\strut#1\strut}}
\newcommand{\CSLIndent}[1]{\hspace{\cslhangindent}#1}



\setlength{\emergencystretch}{3em} % prevent overfull lines

\providecommand{\tightlist}{%
  \setlength{\itemsep}{0pt}\setlength{\parskip}{0pt}}



 


\DeclareMathOperator{\plim}{plim}
\DeclareMathOperator{\prob}{prob}
\usepackage{booktabs}
\usepackage{longtable}
\usepackage{mathtools}
\usepackage{array}
\usepackage{multirow}
\usepackage{wrapfig}
\usepackage{float}
\usepackage{colortbl}
\usepackage{pdflscape}
\usepackage{tabu}
\usepackage{threeparttable}
\newtheorem{assumption}{Assumption}
\usepackage{fancybox, MnSymbol}
\usepackage{booktabs}
\usepackage{longtable}
\usepackage{array}
\usepackage{multirow}
\usepackage{wrapfig}
\usepackage{float}
\usepackage{colortbl}
\usepackage{pdflscape}
\usepackage{tabu}
\usepackage{threeparttable}
\usepackage{threeparttablex}
\usepackage[normalem]{ulem}
\usepackage{makecell}
\usepackage{xcolor}
\makeatletter
\@ifpackageloaded{caption}{}{\usepackage{caption}}
\AtBeginDocument{%
\ifdefined\contentsname
  \renewcommand*\contentsname{Table of contents}
\else
  \newcommand\contentsname{Table of contents}
\fi
\ifdefined\listfigurename
  \renewcommand*\listfigurename{List of Figures}
\else
  \newcommand\listfigurename{List of Figures}
\fi
\ifdefined\listtablename
  \renewcommand*\listtablename{List of Tables}
\else
  \newcommand\listtablename{List of Tables}
\fi
\ifdefined\figurename
  \renewcommand*\figurename{Figure}
\else
  \newcommand\figurename{Figure}
\fi
\ifdefined\tablename
  \renewcommand*\tablename{Table}
\else
  \newcommand\tablename{Table}
\fi
}
\@ifpackageloaded{float}{}{\usepackage{float}}
\floatstyle{ruled}
\@ifundefined{c@chapter}{\newfloat{codelisting}{h}{lop}}{\newfloat{codelisting}{h}{lop}[chapter]}
\floatname{codelisting}{Listing}
\newcommand*\listoflistings{\listof{codelisting}{List of Listings}}
\makeatother
\makeatletter
\makeatother
\makeatletter
\@ifpackageloaded{caption}{}{\usepackage{caption}}
\@ifpackageloaded{subcaption}{}{\usepackage{subcaption}}
\makeatother

\usepackage{bookmark}
\IfFileExists{xurl.sty}{\usepackage{xurl}}{} % add URL line breaks if available
\urlstyle{same}
\hypersetup{
  pdftitle={Categorical Outcome Models},
  pdfauthor={Prof.~Steven F. Koch},
  hidelinks,
  pdfcreator={LaTeX via pandoc}}


\title{\texorpdfstring{Categorical Outcome
Models}{Categorical Outcome Models}}
\subtitle{\texorpdfstring{How and Why Are They
Used?}{How and Why Are They Used?}}
\author{Prof.~Steven F. Koch}
\date{2026-05-04}

\begin{document}
\frame{\titlepage}

\renewcommand*\contentsname{Table of contents}
\begin{frame}[allowframebreaks]
  \frametitle{Table of contents}
  \setcounter{tocdepth}{3}
  \tableofcontents
\end{frame}
\setcounter{tocdepth}{3}
\tableofcontents
}

\subsection{NOTE}\label{note}

\begin{frame}{NOTE}
\begin{enumerate}
\item
  The R code is not identical across the slides
\item
  Sorry, but html and pdf have a few differences that required some
  changes.
\item
  I have tried to make the substantive features the same.
\item
  I have not tried to show all of the code
\item
  But, there are likely a few differences
\end{enumerate}
\end{frame}

\section{Introduction}\label{introduction}

\subsection{Some Background}\label{some-background}

\begin{frame}{Applied research questions}
\protect\phantomsection\label{applied-research-questions}
\begin{enumerate}
\item
  Does education influence employment prospects?
\item
  Does the size of the household reduce access to:

  \begin{enumerate}
  [a.]
  \tightlist
  \item
    food?
  \item
    healthcare?
  \item
    clean energy?
  \end{enumerate}
\item
  Which illnesses lead to catastrophic health care expenditure?
\item
  Does alcohol/drug use/debt influence mental health?
\item
  Does the age of the predidential candidate affect willingness to vote?
\end{enumerate}
\end{frame}

\begin{frame}{What is the common denominator?}
\protect\phantomsection\label{what-is-the-common-denominator}
\begin{itemize}
\item
  The outcome of interest is categorical:

  \begin{itemize}
  \tightlist
  \item
    employment prospects: employed, unemployed, non-participant
  \item
    access to food: inadequate, adequate, more than adequate
  \item
    catastrophic health care expenditure: catastrophic, non-catastrophic
  \item
    mental health: depressed, not depressed
  \item
    voting: cast ballot or not
  \end{itemize}
\item
  Categorical outcomes are not Gaussian
\item
  Errors are synonymous with mis-classification
\end{itemize}
\end{frame}

\begin{frame}{How should we think about modelling?}
\protect\phantomsection\label{how-should-we-think-about-modelling}
\begin{itemize}
\tightlist
\item
  Two options:

  \begin{enumerate}
  \tightlist
  \item
    Assume/Pretend from unobserved continuous variable
  \item
    Use a prediction algorithm
  \item
    Sometimes, these are the same
  \end{enumerate}
\item
  Train (estimate) the model:

  \begin{enumerate}
  \tightlist
  \item
    Use all the data, no real evaluation options
  \item
    Create a training set, evaluate on test set
  \end{enumerate}
\end{itemize}
\end{frame}

\subsection{Some Data}\label{some-data}

\begin{frame}{The data}
\protect\phantomsection\label{the-data}
\begin{enumerate}
\item
  We will use the LCS data (Statistics South Africa 2017)
\item
  We will focus on food adequacy
\item
  Our controls will be:

  \begin{itemize}
  \tightlist
  \item
    Household income
  \item
    Gender of household head
  \item
    Number of children in household
  \item
    Number of adults in household
  \end{itemize}
\item
  Data is limited

  \begin{itemize}
  \tightlist
  \item
    Only households headed by Africans
  \item
    Randomly chosen subset of 2000
  \end{itemize}
\end{enumerate}
\end{frame}

\begin{frame}{Data summary}
\protect\phantomsection\label{data-summary}
\tiny

\begin{table}

\caption{\label{tab:dstat}Household summary statistics by self-assessed food adequacy.}
\centering
\begin{tabular}[t]{l|rrrl|rrrl|rrrl|rrr}
\toprule
 & Below (N = 586) & Adequate (N = 1316) & Above (N = 98)\\
\midrule
\bf{Household head gemder} & ~ & ~ & ~\\
~~ Male headed & 52.39\% & 55.32\% & 70.41\%\\
\bf{Household Composition} & ~ & ~ & ~\\
~~ Children: mean (sd) & 1.21 (1.28) & 1.07 (1.19) & 0.79 (1.06)\\
~~ Adults: mean (sd) & 2.46 (1.33) & 2.40 (1.30) & 2.38 (1.24)\\
\bf{Inverse hyperbolic sine of income} & ~ & ~ & ~\\
~~ IHS: mean (sd) & 8.31 (1.26) & 9.03 (1.17) & 9.67 (1.20)\\
\bottomrule
\end{tabular}
\end{table}

\normalsize

\begin{itemize}
\tightlist
\item
  We will consider binary: below v adequate
\item
  We will consider multinomial: all levels
\end{itemize}
\end{frame}

\subsection{Setup Concepts}\label{setup-concepts}

\begin{frame}{Binary model structure}
\protect\phantomsection\label{binary-model-structure}
\begin{itemize}
\tightlist
\item
  A series of observations \((y_i,X_i)\)
\item
  Machine learning thinks of these as training data
\item
  \(X_i\) is observed characteristics and
\end{itemize}

\[
y_i =  \text{Food adequacy} = \begin{cases}
\text{Below adequate} \\
\text{Adequate}
\end{cases}
\]

\begin{itemize}
\tightlist
\item
  Often, we impose `arbitrary' numerical values, \(\{0, 1 \}\)
\item
  Should be intuitive
\end{itemize}

\[
y_i =  \text{Food adequacy} = \begin{cases}
0 & \text{ if Below adequate} \\
1 & \text{ if Adequate}
\end{cases}
\]
\end{frame}

\begin{frame}{Multinomial ranked responses}
\protect\phantomsection\label{multinomial-ranked-responses}
\begin{itemize}
\tightlist
\item
  Keeping all outcomes
\item
  One could treat them as ranked

  \begin{itemize}
  \tightlist
  \item
    bad to good
  \item
    good to bad
  \end{itemize}
\end{itemize}

\[
y_i = \text{Food adequacy} = \begin{cases}
\text{Below adequate} \\
\text{Adequate} \\
\text{More than adequate}
\end{cases}
\]

\begin{itemize}
\tightlist
\item
  It is often preferred that the spacing is both

  \begin{itemize}
  \tightlist
  \item
    obvious, and
  \item
    constant
  \end{itemize}
\item
  Here, it is obvious, maybe not constant `steps'
\item
  An example with both?
\end{itemize}
\end{frame}

\begin{frame}{Multinomial unranked outcomes}
\protect\phantomsection\label{multinomial-unranked-outcomes}
\begin{itemize}
\tightlist
\item
  If the ordering is not obvious, it is better to not order them
\end{itemize}

\[
y_i = \text{Employment Status} = \begin{cases}
\text{Employed} \\
\text{Unemployed} \\
\text{Non-participant in the labour force}
\end{cases}
\]

\begin{itemize}
\tightlist
\item
  Or
\end{itemize}

\[
y_i = \text{Favourite team} = \begin{cases}
\text{Arsenal} \\
\text{Manchester City} \\
\text{Liverpool} \\
\text{Manchester United}
\end{cases}
\]
\end{frame}

\section{Binary Modelling}\label{binary-modelling}

\subsection{Logit/Probit/LPM}\label{logitprobitlpm}

\begin{frame}{Main ideas}
\protect\phantomsection\label{main-ideas}
\begin{itemize}
\tightlist
\item
  Estimate relationship between \(X_i\) and `success' or `failure'.
\item
  `Success': \(E[y_i|X_i] =\)prob\((y_i=1|X_i)\).
\item
  `Failure': \(E[y_i|X_i] =\)prob\((y_i=0|X_i)\).
\item
  How to model probability?

  \begin{itemize}
  \tightlist
  \item
    Linearly
  \item
    With a probability density function, i.e., MLE
  \item
    Some other way?
  \end{itemize}
\end{itemize}
\end{frame}

\begin{frame}{Assume/pretend continuous}
\protect\phantomsection\label{assumepretend-continuous}
\begin{itemize}
\tightlist
\item
  Noisy measure success \(Z_i\gamma_k+e_{ik}\)
\item
  Outcome determined by larger measure from \(k=\{0,1\}\)
\end{itemize}

\[
\begin{split}
\prob(y_i=1|X_i) &= \prob(X_i\gamma_1+e_{i1}>X_i\gamma_0 +e_{i0}) \\
&= \prob(X_i\gamma_1 - X_i\gamma_0>e_{i0}-e_{i1}) \\
&= \prob(e_{i0}-e_{i1}<X_i(\gamma_1-\gamma_0)) \\
&= \prob(u_i<X_i\beta)
\end{split}
\]

\begin{itemize}
\tightlist
\item
  Make probability distribution assumption

  \begin{itemize}
  \tightlist
  \item
    If we allow \(u\) to be normal, we have a probit
  \item
    If we allow \(u\) to be logistic, we have a logit
  \item
    If we allow the probability to be linear, we have an LPM
  \end{itemize}
\end{itemize}
\end{frame}

\begin{frame}{Linear Probability Model (LPM)}
\protect\phantomsection\label{linear-probability-model-lpm}
\begin{itemize}
\tightlist
\item
  Advantages

  \begin{itemize}
  \tightlist
  \item
    Easy to estimate
  \item
    Easy to interpret
  \end{itemize}
\item
  Disadvantages

  \begin{itemize}
  \tightlist
  \item
    Probability laws may be violated
  \item
    Residual is heteroskedastic (easy to fix, though)
  \item
    Marginal effects potentially unreasonable
  \end{itemize}
\end{itemize}
\end{frame}

\begin{frame}{LPM probability issues?}
\protect\phantomsection\label{lpm-probability-issues}
\begin{figure}[H]

{\centering \includegraphics[width=0.6\linewidth,height=0.6\textheight]{Categoris_Day11_pdfslides_files/figure-beamer/binary-pic-1.pdf}

}

\caption{Illustration of Binary Outcome (jittered) Against Control
Variable Along with Jump Function Guess and a Smooth Function Guess.}

\end{figure}%
\end{frame}

\subsection{Maximum Likelihood Estimation
(MLE)}\label{maximum-likelihood-estimation-mle}

\begin{frame}{Forcing the Probability to be a Probability}
\protect\phantomsection\label{forcing-the-probability-to-be-a-probability}
\begin{itemize}
\tightlist
\item
  Rather than linear probability model
\item
  Assume a probability density - keep to one control for simplicity
\item
  logit:
\end{itemize}

\[
\prob(y_i=1|X_i) = \Lambda(\alpha + X_i\beta) = 
     \frac{\exp(\alpha + X_i\beta)}{1+\exp(\alpha + X_i\beta)}
\]

\begin{itemize}
\tightlist
\item
  probit:
\end{itemize}

\[
\prob(y_i=1|X_i) = \Phi(X_i\beta) = \int_{-\infty}^{X_i\beta} \phi(z)dz
\]
\end{frame}

\begin{frame}{Logit derivative note}
\protect\phantomsection\label{logit-derivative-note}
\begin{itemize}
\tightlist
\item
  Given
\end{itemize}

\[
\Lambda(z) = \frac{e^z}{1+e^z} 
\]

\begin{itemize}
\tightlist
\item
  The derivative is quite simple
\item
  That is why it is still used over probit in most other fields
\end{itemize}

\[
\begin{split}
\frac{\partial \Lambda(z)}{\partial z} &= \frac{e^z}{1+e^z} -
\frac{e^{2z}}{(1+e^z)^2} \\
 &= \frac{(1+e^z)e^z}{(1+e^z)^2}- \frac{e^{2z}}{(1+e^z)^2} \\
 &= \frac{e^z}{(1+e^z)^2} = \frac{e^z}{1+e^z}\frac{1}{1+e^z} \\
 &= \Lambda(z)(1-\Lambda(z))
\end{split}
\]
\end{frame}

\begin{frame}{MLE from Bernoulli}
\protect\phantomsection\label{mle-from-bernoulli}
\begin{itemize}
\tightlist
\item
  Assume a probability mass function for likelihood:
\item
  Recall, \(y_i = 0\) or \(y_i = 1\)
\end{itemize}

\[
f(y_i|X_i) = p_i^{y_i} (1-p_i)^{(1-y_i)}
\]

\begin{itemize}
\tightlist
\item
  Assuming IID and logit yields the likelihood:
\end{itemize}

\[
\mathcal{L} = \prod_{i=1}^N  \Lambda(\alpha + X_i\beta)^{y_i} \big[1-\Lambda(\alpha + X_i\beta)\big]^{(1-y_i)}
\]

\begin{itemize}
\tightlist
\item
  The log-likelihood
\end{itemize}

\[
\ell = \ln \mathcal{L} = \sum_{i=1}^N y_i \cdot \ln \Lambda(\alpha + X_i\beta) + 
   (1-y_i) \cdot \ln \big[1-\Lambda(\alpha + X_i\beta)\big]
\]
\end{frame}

\begin{frame}{Logit MLE score equation for \(\beta\)}
\protect\phantomsection\label{logit-mle-score-equation-for-beta}
\begin{itemize}
\tightlist
\item
  Maximize the log-likelihood over model parameters
\end{itemize}

\[
\begin{split}
\frac{\partial \ell}{\partial \beta} &= \sum_{i=1}^N y_i \frac{\partial
  \Lambda / \partial \beta}{ \Lambda (\alpha + X_i\beta)} + (1-y_i) 
   \left[ \frac{- \partial  \Lambda / \partial \beta}{1-\Lambda(\alpha + X_i\beta)} \right] = 0 \\
&= \sum_{i=1}^N y_i X_i \frac{\big[\Lambda(1-\Lambda)\big]}{\Lambda} 
  + (1-y_i)(-X_i) \frac{\big[\Lambda(1-\Lambda)\big]}{1-\Lambda}  = 0 \\
&= \sum_{i=1}^N X_i \big[y_i - \Lambda(\alpha + X_i \beta) \big] = 0
\end{split}
\]
\end{frame}

\begin{frame}{Logit MLE score equation for \(\alpha\)}
\protect\phantomsection\label{logit-mle-score-equation-for-alpha}
\[
\begin{split}
\frac{\partial \ell}{\partial \alpha} &= \sum_{i=1}^N y_i \frac{\partial
  \Lambda / \partial \alpha}{ \Lambda (\alpha + X_i\beta)} + (1-y_i) 
   \left[ \frac{- \partial  \Lambda / \partial \alpha}{1-\Lambda(\alpha + X_i\beta)} \right] = 0 \\
&= \sum_{i=1}^N y_i \frac{\big[\Lambda(1-\Lambda)\big]}{\Lambda} 
  + (1-y_i) \frac{\big[\Lambda(1-\Lambda)\big]}{1-\Lambda}  = 0 \\
&= \sum_{i=1}^N  \big[ y_i - \Lambda(\alpha + X_i \beta) \big] = 0 \\
\hat{\Lambda} &= \overline{y}
\end{split}
\]

\begin{itemize}
\tightlist
\item
  Must solve both (all) score equations simultaneously
\item
  Estimation based on numerical methods
\item
  Thus, slower than linear model, too.
\end{itemize}
\end{frame}

\begin{frame}{Estimates of interest from MLE?}
\protect\phantomsection\label{estimates-of-interest-from-mle}
\begin{itemize}
\tightlist
\item
  What does a parameter tell us?
\item
  Recall the logit probability
\end{itemize}

\[
\prob(y_i=1|X_i) = \Lambda(\alpha + X_i\beta)
\]

\begin{itemize}
\tightlist
\item
  The marginal effect of \(X_i\):
\end{itemize}

\[
\begin{split}
\frac{\partial \prob(y_i=1|X_i)}{\partial X_i} &= \beta
\frac{\partial \Lambda}{\partial X_i} \\
 &= \beta \Lambda (1-\Lambda) 
\end{split}
\]
\end{frame}

\begin{frame}{Interpretation}
\protect\phantomsection\label{interpretation}
\begin{itemize}
\item
  From the estimates, we see

  \begin{itemize}
  \tightlist
  \item
    \(\beta\) has the same sign as the marginal effect
  \item
    BUT, IT IS NOT THE MARGINAL EFFECT
  \item
    So, limited value to parameter estimates
  \end{itemize}
\item
  Why might LPMs be preferred?
\item
  Final comments

  \begin{itemize}
  \tightlist
  \item
    I am assuming continuity of \(X\)
  \item
    If not continuous, then the derivative is a ``difference''
  \item
    Probit is more complicated math
  \item
    Same qualitative conclusion
  \end{itemize}
\item
  We will return to this
\end{itemize}
\end{frame}

\subsection{Variances}\label{variances}

\begin{frame}{Basic MLE Variance}
\protect\phantomsection\label{basic-mle-variance}
\begin{itemize}
\tightlist
\item
  Variances

  \begin{itemize}
  \tightlist
  \item
    Outerproduct of the score
  \item
    Derivative of the
    score.\footnote{The mathematics is greatly simplified by noting that $E[y_i-F]=0$.}
  \end{itemize}
\end{itemize}

\[
V[\beta] = \left( \sum_{i=1}^N \frac{[(\partial F/ \partial \beta)^2]
    X_i X_i'}{F(1-F)} \right)^{-1}
\]

\begin{itemize}
\tightlist
\item
  Use estimates for actual values (inference is `simple'):
\end{itemize}

\[
\hat{\beta} \xrightarrow{d} \mathcal{N}(\beta, \mathcal{I}(\hat{\beta})^{-1})
\]

\begin{itemize}
\tightlist
\item
  But, inference on parameters may not mean much\ldots{}
\end{itemize}
\end{frame}

\begin{frame}{Second derivatives under logit}
\protect\phantomsection\label{second-derivatives-under-logit}
\[
\begin{split}
\frac{\partial \ell}{\partial \alpha} &=  \sum_{i=1}^N  \big[ y_i - \Lambda(\alpha + X_i \beta) \big] \\
\frac{\partial^2 \ell}{\partial \alpha^2} &= - \sum_{i=1}^N \Lambda(1-\Lambda)
\end{split}
\]

\[
\begin{split}
\frac{\partial \ell}{\partial \beta} &= \sum_{i=1}^N X_i \big[y_i - \Lambda(\alpha + X_i \beta) \big]  \\
\frac{\partial^2 \ell}{\partial \beta^2} &= - \sum_{i=1}^N X_i^2 \cdot \Lambda(1-\Lambda)
\end{split}
\]
\end{frame}

\begin{frame}{One more derivative and the variance matrix}
\protect\phantomsection\label{one-more-derivative-and-the-variance-matrix}
\[
\begin{split}
\frac{\partial \ell}{\partial \beta} &= \sum_{i=1}^N X_i \big[y_i - \Lambda(\alpha + X_i \beta) \big]  \\
\frac{\partial^2 \ell}{\partial \beta \partial \alpha} &= - \sum_{i=1}^N X_i \cdot \Lambda(1-\Lambda)
\end{split}
\]

\begin{itemize}
\tightlist
\item
  Matrix \(\Omega = [\alpha \text{ } \beta]'\)
\end{itemize}

\[
V[\hat\Omega] = \begin{bmatrix}
- \sum_{i=1}^N \Lambda(1-\Lambda) &  - \sum_{i=1}^N X_i \cdot \Lambda(1-\Lambda) \\
- \sum_{i=1}^N X_i \cdot \Lambda(1-\Lambda) & - \sum_{i=1}^N X_i^2 \cdot \Lambda(1-\Lambda)
\end{bmatrix}^{-1}
\]
\end{frame}

\begin{frame}{Delta method}
\protect\phantomsection\label{delta-method}
\begin{itemize}
\tightlist
\item
  Remember, our interest is in the marginal effect
\item
  So, we need its variance, not the one for the parameters
\item
  The basic point:

  \begin{itemize}
  \tightlist
  \item
    From MLE:
    \(\sqrt{n}(\Omega - \hat\Omega) \rightarrow \mathcal{N}(0,\Sigma)\)
  \item
    Thus
  \end{itemize}
\end{itemize}

\[
\sqrt{n}\left(g(\Omega) - g\left(\hat\Omega \right) \right) \rightarrow \mathcal{N}\left(0,\nabla g(\Omega)' \cdot \Sigma \cdot \nabla g(\Omega)\right)
\]

\begin{itemize}
\tightlist
\item
  This is similar to the variance rule for squaring a constant
\item
  The proof of this relies on Taylor expansions
\end{itemize}
\end{frame}

\begin{frame}{Marginal effect variance}
\protect\phantomsection\label{marginal-effect-variance}
\begin{itemize}
\tightlist
\item
  Given the marginal effect of \(X_i\):
\end{itemize}

\[
\frac{\partial \prob(y_i=1|X_i)}{\partial X_i} = \beta \Lambda (1-\Lambda) 
\]

\begin{itemize}
\tightlist
\item
  We need to take on the derivative of this (squaring it)
\item
  But, it is a three part product
\end{itemize}

\[
\begin{split}
\frac{\partial^2 \prob(y_i=1|X_i)}{\partial X_i \partial \beta} &= \beta \Lambda (1-\Lambda) \\
 &= \Lambda (1-\Lambda) + X_i \beta \Lambda (1-\Lambda)^2 - X_i\beta\Lambda^2(1-\Lambda)
 \end{split}
\]

\begin{itemize}
\tightlist
\item
  Or, we can bootstrap..
\item
  Depends a bit on what we want to consider
\end{itemize}
\end{frame}

\section{Other Approaches}\label{other-approaches}

\subsection{Background}\label{background}

\begin{frame}[allowframebreaks]{Some other estimation concepts}
\protect\phantomsection\label{some-other-estimation-concepts}
\begin{itemize}
\item
  One can use many other probability functions\ldots{}
\item
  Other distribution functions can be used

  \begin{itemize}
  \tightlist
  \item
    Normality - probit
  \item
    Complementary log-log: \(C(X_i\beta) = 1-\exp(-exp(X_i\beta))\).
  \item
    Or \(\beta\) or \(\Gamma\) distributions
  \item
    These are not commonly applied
  \end{itemize}
\item
  Discriminant analysis

  \begin{itemize}
  \tightlist
  \item
    Linear DA
  \item
    Quadratic DA
  \item
    We will discuss this
  \end{itemize}
\item
  Maximum Score Estimation (Manski 1975)

  \begin{itemize}
  \tightlist
  \item
    Estimate the median
    \(Q_N(\beta) = \sum_i |y_i-\mathbf{1}(X_i\beta>0)|\)
  \item
    Quantile regression beyond the median (Koenker and Bassett 1978)
  \end{itemize}
\item
  Semiparametric Methods - do not specify \(F\).

  \begin{itemize}
  \tightlist
  \item
    Klein and Spady (1993) is an early approach
  \item
    There are others
  \end{itemize}
\item
  Fully Nonparametric Methods

  \begin{itemize}
  \tightlist
  \item
    Similar to what we discussed with densities and distributions
  \item
    Also, there are clustering approaches
  \end{itemize}
\end{itemize}
\end{frame}

\begin{frame}{Correct Classification Ratio}
\protect\phantomsection\label{correct-classification-ratio}
\begin{itemize}
\tightlist
\item
  Map observed outcomes \(Y_i=k\), \(j \times 1\) vector \(\Upsilon_i\)
\end{itemize}

\[
\Upsilon_{ik} = \begin{cases}
 1 & \text{if } Y_i=k \\
 0 & \text{otherwise}
\end{cases}
\]

\begin{itemize}
\tightlist
\item
  Map predicted outcomes similarly: \(\hat{\Upsilon}_i\)
\end{itemize}

\[
\hat{\Upsilon}_{ik} = \begin{cases}
 1 & \text{if } \widehat{p}_{ik} = \max_j \{\widehat{p}_{ij} \} \\
 0 & \text{otherwise}
\end{cases}
\]
\end{frame}

\begin{frame}{Correct Classification Ratio II}
\protect\phantomsection\label{correct-classification-ratio-ii}
\begin{itemize}
\tightlist
\item
  Define a classification function
\item
  Penalize incorrect predictions
\end{itemize}

\[
Q_i(\Upsilon,\hat{\Upsilon},n) = \begin{cases}
 1 & \text{if } \Upsilon_i = \hat{\Upsilon}_i \\
 0 & \text{otherwise}
\end{cases}
\]

\begin{itemize}
\tightlist
\item
  Leads to correct classification ratio (CCR)
\end{itemize}

\[
CCR = n^{-1} \sum_{i=1}^n Q_i(\Upsilon,\hat{\Upsilon},n) 
\]

\begin{itemize}
\tightlist
\item
  Develop a simple table
\item
  Actual against predicted outcomes
\end{itemize}

\[
CM = \Upsilon' \hat{\Upsilon}
\]
\end{frame}

\subsection{Approaches to Bayes'
classificaion}\label{approaches-to-bayes-classificaion}

\begin{frame}{Bayes' classification: what is it?}
\protect\phantomsection\label{bayes-classification-what-is-it}
\begin{enumerate}
\tightlist
\item
  Focus is a bit more on prediction
\item
  Still something like before:

  \begin{itemize}
  \tightlist
  \item
    A series of observations \((y_i,X_i)\)
  \item
    Machine learning thinks of these as training data
  \item
    But, we keep outcomes qualitative
  \end{itemize}
\end{enumerate}

\[
y_i =  \text{Food adequacy} = \begin{cases}
\text{Below adequate} \\
\text{Adequate}
\end{cases}
\]

\begin{enumerate}
\setcounter{enumi}{1}
\tightlist
\item
  Our interest -- choose max value

  \begin{itemize}
  \tightlist
  \item
    With two classes, must exceed 0.5
  \item
    With three classes, not as simple
  \end{itemize}
\end{enumerate}

\[
\prob(y_i = j | X_i = x_0)
\]
\end{frame}

\begin{frame}{We will generalise this a bit}
\protect\phantomsection\label{we-will-generalise-this-a-bit}
\begin{itemize}
\tightlist
\item
  Still using same basic probability model
\item
  Consider a specific outcome \(k\) from \(J\) options
\item
  Technically, we could drop one outcome as a baseline
\end{itemize}

\[
\prob(y_i = j | X_i = x_0) = \frac{\pi_k f_k(x)} {\sum_{\ell \in J} \pi_\ell f_\ell(x)}
\]

\begin{itemize}
\tightlist
\item
  How do we find \(\pi\) and \(f\)?

  \begin{itemize}
  \tightlist
  \item
    \(\pi\) can be pulled from the data
  \item
    \(f\) is often assumed
  \item
    Note: \(f\) is over the control/pedictor variables
  \end{itemize}
\end{itemize}
\end{frame}

\begin{frame}{\(k-\) Nearest Neighbors}
\protect\phantomsection\label{k--nearest-neighbors}
\begin{enumerate}
\tightlist
\item
  One way to operationalise this classifier

  \begin{itemize}
  \tightlist
  \item
    Estimate \(\pi\) from subsets of the data
  \item
    \(f\) is then about choosing subsets
  \end{itemize}
\item
  Basic intuition:

  \begin{itemize}
  \tightlist
  \item
    Assign class based on similar observations
  \item
    This depends on how ``wide'' we make the comparison group
  \item
    Mathematically
  \end{itemize}
\end{enumerate}

\[
\prob(y_i = j | X_i = x_0) = \frac{1}{k} \sum_{\ell \in N_k} \mathbf{1}(y_\ell = j)
\]
\end{frame}

\begin{frame}{Explaining process for \(k-\)NN}
\protect\phantomsection\label{explaining-process-for-k-nn}
\begin{enumerate}
\tightlist
\item
  Choose \(k\)
\item
  Find \(k\) neighbors around an observation
\item
  We see how they are classified, on average
\item
  We choose largest average
\item
  \(k\) is like a bandwidth

  \begin{itemize}
  \tightlist
  \item
    Estimates are sensitive to size of \(k\)
  \item
    Need a way to `optimise'
  \item
    We discuss cross-validation ideas later in semester
  \end{itemize}
\end{enumerate}
\end{frame}

\begin{frame}{LDA}
\protect\phantomsection\label{lda}
\begin{itemize}
\tightlist
\item
  Instead, assume \(f\) comes from normal
\item
  Implicitly, \(x\) is quantitative (but does not have to be)
\item
  Estimate \(\pi\) from the training data
\item
  Again, we might need to think about cross-validation
\item
  Below, only one \(x\), easily generalised\ldots{}
\end{itemize}

\[
f_k(x) = \frac{1}{\sigma_k \sqrt{2\pi}} \cdot \exp \left[ -\frac{1}{2\sigma^2_k} \left(x - \mu_k \right)^2 \right]
\]
\end{frame}

\begin{frame}{Linear discriminant choice}
\protect\phantomsection\label{linear-discriminant-choice}
\begin{itemize}
\tightlist
\item
  Assume \(J=2\) and constant variance across classes
\item
  Take logs of the basic probability relationship
\item
  Yields the basic linear result
\item
  Find the largest value of
\end{itemize}

\[
\delta_k(x) = x \cdot \frac{\mu_k}{\sigma^2} - \frac{\mu_k^2}{2\sigma^2}+\ln \pi_k
\]

\begin{itemize}
\tightlist
\item
  It is linear in \(x\), so it is linear discriminant analysis
\item
  Simply put sample estimates in for population estimates
\item
  In principle, one could assume different forms for \(f\)
\item
  Certainly, one can allow for more classes
\end{itemize}
\end{frame}

\begin{frame}{Quadratic discriminant analysis}
\protect\phantomsection\label{quadratic-discriminant-analysis}
\begin{itemize}
\tightlist
\item
  Allows for different variances in different classes
\item
  Variances include squared terms
\item
  Thus, the discriminator includes a quadratic term in \(x\)
\end{itemize}

\[
\delta(x) = -\frac{1}{2} \left(x-\mu_k \right)' \Sigma^{-1}_k \left(x-\mu_k \right) -\frac{1}{2} \ln \left| \Sigma_k \right|  + \ln \pi 
\]

\begin{itemize}
\tightlist
\item
  The first term includes \(x'x\)
\item
  It is squared terms and interaction terms
\item
  Technically, variances and covariances
\end{itemize}
\end{frame}

\begin{frame}{Naïve Bayes'}
\protect\phantomsection\label{nauxefve-bayes}
\begin{itemize}
\tightlist
\item
  In LDA and QDA, we have a joint normal
\item
  That means the distributions are `correlated'
\item
  Instead, assume independence

  \begin{itemize}
  \tightlist
  \item
    e.g., a series of independent normals
  \item
    Basically, QDA with diagonal \(\Sigma\)
  \end{itemize}
\item
  Thus, we replace \(f_k\) with
  \(f_{k1}(x_1) \times f_{k2}(x_2) \times \cdots\)
\end{itemize}

\[
\prob(y_i = j | X_i = x_0) = \frac{\pi_k \times f_{k1}(x_1) \times f_{k2}(x_2) \times \cdots \times f_{kJ}(x_J)} {\sum_{\ell \in J} \pi_\ell \times f_{\ell1}(x_1) \times f_{\ell2}(x_2) \times \cdots \times f_{\ell J}(x_J)}
\]

\begin{itemize}
\tightlist
\item
  We need \(\pi_k\) and densities for each \(x\)

  \begin{itemize}
  \tightlist
  \item
    \(\pi_k\) as average outcomes in training data
  \item
    \(x\) can be estimated from histogram or np density
  \end{itemize}
\end{itemize}
\end{frame}

\subsection{Predictive performance}\label{predictive-performance}

\begin{frame}{Pseudo-\(R^2\) and Other Performance Measures}
\protect\phantomsection\label{pseudo-r2-and-other-performance-measures}
\begin{itemize}
\tightlist
\item
  A simple performance measure
\item
  Like OLS, use the intercept model as baseline
\end{itemize}

\[
\begin{split}
R^2_{b} &=
1-\frac{ \mathcal{L}_N ( \hat{ \beta} )} {\mathcal{L}_N (\overline{y})} \\
 &= 1 - \frac{ \sum_{i=1}^N y_i \ln \hat{p}_i + (1-y_i) \ln
 (1-\hat{p}_i)} {N [ \overline{y} \ln \overline{y} +(1- \overline{y})
 \ln (1-\overline{y} )] }
\end{split}
\]

\begin{itemize}
\tightlist
\item
  Other Comparisons

  \begin{itemize}
  \tightlist
  \item
    Hosmer-Lemeshow
  \item
    ROC Curve and Classification Matrix
  \end{itemize}
\item
  We should maybe pay attention to outliers, as well
\end{itemize}
\end{frame}

\begin{frame}{Hosmer-Lemeshow}
\protect\phantomsection\label{hosmer-lemeshow}
\begin{itemize}
\tightlist
\item
  It is a Pearson \(\chi^2\) test
\end{itemize}

\[
\chi_{G-2}^2 = \sum_{g=1}^G \frac{(\text{observed}_g -
  \text{expected}_g)^2}{\text{expected}_g}
\]

\begin{itemize}
\tightlist
\item
  Observed successes/failures in the data
\item
  \(G\) subgroups of analysis

  \begin{itemize}
  \tightlist
  \item
    2 required, more ok
  \item
    Based on the predicted ``probabilities''classes''
  \end{itemize}
\item
  Expected successes/failures

  \begin{itemize}
  \tightlist
  \item
    predicted in the model
  \item
    Assume phat\$\textgreater\$0.5 as threshold
  \end{itemize}
\item
  \emph{Good model works in subgroups, too}
\end{itemize}
\end{frame}

\begin{frame}{ROC and Sensitivity}
\protect\phantomsection\label{roc-and-sensitivity}
Receiver operating characteristic (ROC) curve plots binary
classification as threshold changes. It gives us some idea of
sensitivity of results to arbitrary choice of threshold. The area under
the curve tells us the probability that a classifier will rank a
randomly chosen success higher than a randomly chosen failure, and is
similar in spirit to Mann-Whitney and Wilcoxon rank tests.

\vspace{3mm}

\begin{itemize}
\tightlist
\item
  Classifications depend on the chosen cutoff
\item
  Analysis of tradeoff between sensitivity and specificity
\item
  ROC curve shows whether model does better than a ``guess''
\item
  Available in STATA and R, and others
\item
  Not commonly used in economics, though
\end{itemize}
\end{frame}

\begin{frame}{Brier Score}
\protect\phantomsection\label{brier-score}
\begin{itemize}
\tightlist
\item
  Rather than using probabilities to get ``classes''
\item
  USE the probabilities
\end{itemize}

\[
\text{Brier} = N^{-1} \sum_{i=1}^N \left( \hat{p}_i - o_i  \right)^2
\]

\begin{itemize}
\tightlist
\item
  Basically, a mean squared error
\item
  Sort of similar to Hosmer-Lemeshow
\item
  Argued to be a better measure, but not a ``test''
\item
  Lower is better (like AIC/BIC)
\end{itemize}
\end{frame}

\subsection{Binary Example}\label{binary-example}

\begin{frame}{Logit, probit and LPM parameters}
\protect\phantomsection\label{logit-probit-and-lpm-parameters}
\tiny

\begin{table}[!htbp] \centering 
  \caption{} 
  \label{} 
\begin{tabular}{@{\extracolsep{5pt}}lccc} 
\\[-1.8ex]\hline 
\hline \\[-1.8ex] 
 & \multicolumn{3}{c}{\textit{Dependent variable:}} \\ 
\cline{2-4} 
\\[-1.8ex] & bin.adeq & \multicolumn{2}{c}{adequate} \\ 
\\[-1.8ex] & \textit{OLS} & \textit{logistic} & \textit{probit} \\ 
\\[-1.8ex] & (1) & (2) & (3)\\ 
\hline \\[-1.8ex] 
 y.ihs & 0.109$^{***}$ & 0.587$^{***}$ & 0.333$^{***}$ \\ 
  & (0.009) & (0.056) & (0.032) \\ 
  & & & \\ 
 adults & $-$0.028$^{***}$ & $-$0.152$^{***}$ & $-$0.088$^{***}$ \\ 
  & (0.010) & (0.048) & (0.029) \\ 
  & & & \\ 
 kids & $-$0.023$^{**}$ & $-$0.115$^{**}$ & $-$0.068$^{**}$ \\ 
  & (0.010) & (0.050) & (0.030) \\ 
  & & & \\ 
 hhh.maleFemale & 0.012 & 0.049 & 0.027 \\ 
  & (0.022) & (0.114) & (0.068) \\ 
  & & & \\ 
 Constant & $-$0.179$^{**}$ & $-$3.823$^{***}$ & $-$2.127$^{***}$ \\ 
  & (0.079) & (0.461) & (0.267) \\ 
  & & & \\ 
\hline \\[-1.8ex] 
\hline 
\hline \\[-1.8ex] 
\textit{Note:}  & \multicolumn{3}{r}{$^{*}$p$<$0.1; $^{**}$p$<$0.05; $^{***}$p$<$0.01} \\ 
\end{tabular} 
\end{table}

\normalsize
\end{frame}

\begin{frame}{Brief discussion}
\protect\phantomsection\label{brief-discussion}
\begin{enumerate}
\tightlist
\item
  Note: I used a training set of 1647 obs
\item
  All coefficients are different
\item
  Most are statistically significant
\item
  LPM gives average marginal effect

  \begin{itemize}
  \tightlist
  \item
    probit/logit does not
  \item
    probit/logit give direction
  \end{itemize}
\end{enumerate}
\end{frame}

\begin{frame}{Model predictions}
\protect\phantomsection\label{model-predictions}
\tiny

\begin{table}

\caption{\label{tab:binpredict}Predictive performance by model}
\centering
\begin{tabular}[t]{l|r|r|r|r|r|r}
\hline
\multicolumn{1}{c|}{ } & \multicolumn{2}{c|}{LPM} & \multicolumn{2}{c|}{logit} & \multicolumn{2}{c}{probit} \\
\cline{2-3} \cline{4-5} \cline{6-7}
  & Adequate & Below & Adequate & Below & Adequate & Below\\
\hline
Below & 63 & 9 & 63 & 9 & 63 & 9\\
\hline
Adequate & 175 & 3 & 170 & 8 & 173 & 5\\
\hline
\end{tabular}
\end{table}

\normalsize

\begin{itemize}
\tightlist
\item
  All the same for ``Below''
\item
  LPM does better for ``Adequate''
\item
  LPM better overall
\item
  LPM share outside unit interval - 0.008
\end{itemize}
\end{frame}

\begin{frame}{LDA and QDA}
\protect\phantomsection\label{lda-and-qda}
\small

\begin{table}

\caption{\label{tab:discriminant}Predictive performance by model}
\centering
\begin{tabular}[t]{l|r|r|r|r}
\hline
\multicolumn{1}{c|}{ } & \multicolumn{2}{c|}{LDA} & \multicolumn{2}{c}{QDA} \\
\cline{2-3} \cline{4-5}
  & Below & Adequate & Below & Adequate\\
\hline
Below & 9 & 63 & 11 & 61\\
\hline
Adequate & 5 & 173 & 13 & 165\\
\hline
\end{tabular}
\end{table}

\normalsize

\begin{itemize}
\tightlist
\item
  LDA does better for ``Adequate'' and overall
\item
  QDA does better for ``Below''
\item
  LPM still in the lead\ldots{}
\end{itemize}
\end{frame}

\begin{frame}{\(k-\)NN and naïve Bayes'}
\protect\phantomsection\label{k-nn-and-nauxefve-bayes}
\small

\begin{table}

\caption{\label{tab:bayes-knn}Predictive performance by model}
\centering
\begin{tabular}[t]{l|r|r|r|r|r|r}
\hline
\multicolumn{1}{c|}{ } & \multicolumn{2}{c|}{Bayes} & \multicolumn{2}{c|}{KNN 3} & \multicolumn{2}{c}{KNN - 5} \\
\cline{2-3} \cline{4-5} \cline{6-7}
  & Below & Adequate & Below & Adequate & Below & Adequate\\
\hline
Below & 10 & 62 & 20 & 52 & 17 & 55\\
\hline
Adequate & 7 & 171 & 40 & 138 & 26 & 152\\
\hline
\end{tabular}
\end{table}

\normalsize

\begin{itemize}
\tightlist
\item
  \(k-\)NN with 3 neighbours does better for ``Below''
\item
  Bayes does better for ``Adequat'' and overall
\item
  LPM still better overall
\item
  Unusual, but not impossible
\end{itemize}
\end{frame}

\begin{frame}{General \(k-\)NN plot: LPM still better}
\protect\phantomsection\label{general-k-nn-plot-lpm-still-better}
\begin{figure}[H]

{\centering \includegraphics[width=0.6\linewidth,height=0.6\textheight]{Categoris_Day11_pdfslides_files/figure-beamer/knn-plot-1.pdf}

}

\caption{Nearest neighbor performance relative to LPM}

\end{figure}%
\end{frame}

\section{References}\label{references}

\begin{frame}{References \{allowframebreaks\}}
\protect\phantomsection\label{references-allowframebreaks}
\protect\phantomsection\label{refs}
\begin{CSLReferences}{1}{1}
\bibitem[\citeproctext]{ref-kleinspady1993}
Klein, R W, and R H Spady. 1993. {``An Efficient Semiparametric
Estimator of Binary Response Models.''} \emph{Econometrica} 61 (2):
387--421.

\bibitem[\citeproctext]{ref-koenkerbasset1978}
Koenker, Roger, and G Bassett. 1978. {``Regression Quantiles.''}
\emph{Econometrica} 50 (1): 43--61.

\bibitem[\citeproctext]{ref-manski1975}
Manski, C F. 1975. {``Maximum Score Estimation of the Stochastic Utility
Model.''} \emph{Journal of Econometrics} 3 (3): 205--28.

\bibitem[\citeproctext]{ref-lcs_2014_2015}
Statistics South Africa. 2017. {``Living Conditions Survey 2014--2015,
South Africa.''} Version v1.1. DataFirst, University of Cape Town;
Statistics South Africa.
\url{https://microdata.worldbank.org/index.php/catalog/2882}.

\end{CSLReferences}
\end{frame}




\end{document}
